%=============================================================================
\section{多回路架构设计}
%=============================================================================

上一章我们设计了单个 PID 控制器。但真正的比赛机器人不会只跑一个控制回路。云台有电流环、速度环、位置环。平衡车有倾角环和速度环。底盘有每个轮子的速度环，还有上层的轨迹跟踪环。这一章教你怎么把这些回路堆叠成一个协调一致的架构。

一个机器人"凑合能用"和"稳赢比赛"之间的差距，往往就在架构上。两支队伍可以用一模一样的 PID 算法、一模一样的电机、一模一样的传感器——一个机器人稳如磐石，另一个抖个不停。赢的那支队伍搞明白了回路之间怎么配合。这就是我们这章要讲的。

我们先从全局讲起——从传感器到执行器的完整控制栈——然后深入串级控制和让它生效的带宽分离原则。之后讨论任务空间 vs 关节空间控制、上层状态机、以及保证一切准时运行的时序架构。最后用一个完整的云台案例把所有东西串起来。

%-----------------------------------------------------------------------------
\subsection{完整的控制栈}
%-----------------------------------------------------------------------------

每一台比赛机器人——不管你意识到没有——都实现了一个分层流水线：

\[
\text{传感器} \;\longrightarrow\; \text{感知} \;\longrightarrow\; \text{规划} \;\longrightarrow\; \text{控制} \;\longrightarrow\; \text{执行}
\]

\textbf{传感器}产生原始数据：编码器计数、IMU 加速度、电流检测 ADC 值。\textbf{感知}把原始数据变成有用的状态估计：EKF 融合 IMU 和编码器得到倾角，或者视觉流水线提取目标坐标。\textbf{规划}决定机器人下一步该做什么：跟踪轨迹、瞄准目标、保持平衡。\textbf{控制}计算让规划落地的指令：PID 输出、力矩命令、PWM 占空比。\textbf{执行}把指令作用到物理世界：电机驱动切 H 桥、电磁阀触发。

拿一个比赛云台来说。传感器层以 8~kHz 读六轴 IMU、以 10~kHz 读电机编码器。感知层以 1~kHz 跑 EKF 估计偏航和俯仰角。规划层以 120~Hz 从视觉协处理器接收目标坐标，计算期望云台角度。控制层以不同速率跑串级 PID（位置 $\to$ 速度 $\to$ 电流）。执行层是以 20~kHz 开关频率工作的三相电机驱动。

关键洞察：\textit{每一层都是下一层的客户}。规划层不知道也不关心电机驱动怎么工作——它只说"去 30 度"。位置 PID 不知道也不关心磁场定向控制——它只输出一个速度指令。这种关注点分离，才让复杂机器人变得可管理。

每一层跑在哪里？一个典型的比赛机器人：

\begin{center}
\begin{tabular}{l l l}
\hline
\textbf{层} & \textbf{典型频率} & \textbf{硬件} \\
\hline
电流控制 & 10--20~kHz & 电机驱动 MCU（如 C610/C620 电调） \\
速度/位置控制 & 500~Hz--1~kHz & 主控 MCU（STM32F4/H7） \\
感知（EKF） & 500~Hz--1~kHz & 主控 MCU \\
规划/轨迹 & 50--200~Hz & 主控 MCU 或协处理器 \\
视觉/检测 & 30--120~Hz & 协处理器（Jetson、RPi 等） \\
\hline
\end{tabular}
\end{center}

注意规律：越底层跑得越快，越上层跑得越慢。这不是巧合——这是物理规律的直接结果，下面马上讲。

%-----------------------------------------------------------------------------
\subsection{串级控制架构}
%-----------------------------------------------------------------------------

\subsubsection{电流 $\to$ 速度 $\to$ 位置}

三环串级是无刷电机控制的主力架构。让我们从物理上理解\textit{为什么}回路要按这个顺序嵌套——以直流电机为例分析（对 FOC 控制的无刷电机做适当坐标变换后结论相同）。

电气动态方程：
\begin{equation}
L \frac{di}{dt} = V - Ri - K_e \omega
\end{equation}
其中 $L$ 是绕组电感，$R$ 是电阻，$K_e$ 是反电动势常数，$\omega$ 是角速度。电气时间常数 $\tau_e = L/R$，小电机一般 0.1--1~ms。这是系统中\textit{最快}的动态。

机械动态方程：
\begin{equation}
J \frac{d\omega}{dt} = K_t i - b\omega - \tau_{\text{load}}
\end{equation}
其中 $J$ 是转子惯量，$K_t$ 是力矩常数，$b$ 是粘性摩擦，$\tau_{\text{load}}$ 是外部负载。机械时间常数 $\tau_m = J/b$ 一般 10--100~ms——比电气动态慢得多。

最后是运动学关系：
\begin{equation}
\frac{d\theta}{dt} = \omega
\end{equation}

这是纯积分，根本没有自然时间常数——它是最"慢"的动态（积分器的增益从直流开始就以 $-20$~dB/十倍频 的速率下降）。

串级结构就是和物理对齐的：\textbf{最内环}控制\textbf{最快}的动态（电流），\textbf{中间环}控制次快的（速度），\textbf{最外环}控制最慢的（位置）。典型带宽：

\begin{center}
\begin{tabular}{l l l}
\hline
\textbf{回路} & \textbf{带宽} & \textbf{控制器类型} \\
\hline
电流 & 5--10~kHz & PI（不需要 D——被控对象是一阶） \\
速度 & 200~Hz--1~kHz & PID 或 PI \\
位置 & 50--200~Hz & PD 或 PID \\
\hline
\end{tabular}
\end{center}

嵌套框图长这样：位置控制器输出速度设定值，速度控制器输出电流（力矩）指令。对于 CAN 总线电调（C610/C620），这个电流指令通过 CAN 发给电调，电调内部处理电流环和 FOC 换相。对于简单驱动器，电流控制器输出电压（PWM 占空比）。无论哪种方式，每个内环都"包裹"住自己的被控对象动态，在外环看到之前就把它驯服了。

\begin{equation}
r_\theta \;\xrightarrow{C_\theta}\; r_\omega \;\xrightarrow{C_\omega}\; r_i \;\xrightarrow{C_i}\; V \;\xrightarrow{\text{motor}}\; i,\;\omega,\;\theta
\end{equation}

\subsubsection{带宽分离原则}

串级架构之所以管用，靠的是一个强力简化：如果内环比外环\textit{快得多}，外环就可以假装内环不存在。

设内环闭环传递函数为：
\begin{equation}
T_{\text{inner}}(s) = \frac{C_{\text{inner}} G_{\text{inner}}}{1 + C_{\text{inner}} G_{\text{inner}}}
\end{equation}

如果内环带宽 $\omega_{\text{inner}}$ 远大于外环带宽 $\omega_{\text{outer}}$，那么在外环关心的所有频率（$\omega \ll \omega_{\text{inner}}$）上，$|T_{\text{inner}}(j\omega)| \approx 1$ 且 $\angle T_{\text{inner}}(j\omega) \approx 0$。换言之，从外环的视角看 $T_{\text{inner}}(s) \approx 1$——内环就像一个完美的、瞬时响应的执行器。

推导很直接。对于一个调好的内环，闭环近似为一阶系统：
\begin{equation}
T_{\text{inner}}(s) \approx \frac{1}{1 + s/\omega_{\text{inner}}}
\end{equation}

在频率 $\omega \ll \omega_{\text{inner}}$ 处，幅值为：
\[
|T_{\text{inner}}(j\omega)| = \frac{1}{\sqrt{1 + (\omega/\omega_{\text{inner}})^2}} \approx 1 - \frac{1}{2}\left(\frac{\omega}{\omega_{\text{inner}}}\right)^2
\]

相位为 $\angle T_{\text{inner}} \approx -\arctan(\omega/\omega_{\text{inner}}) \approx -\omega/\omega_{\text{inner}}$ 弧度。要让近似的幅值误差小于 5\%、相位滞后小于 10\textdegree{}，大致需要 $\omega_{\text{outer}} < \omega_{\text{inner}} / 5$。由此得出：

\textbf{经验法则：}内环应该比外环快 5--10 倍。

如果违反了这个分离会怎样？假设位置环带宽 200~Hz，速度环带宽 300~Hz——比值只有 1.5 倍。位置环这时候就能"看到"速度环带来的显著相位滞后。在位置环的穿越频率处，速度环贡献了也许 30--40\textdegree{} 的额外相位滞后。这吃掉了位置环的相位裕度，导致超调甚至振荡。两个环\textit{耦合}了——调一个影响另一个，漂亮的分离故事就讲不通了。

这是比赛中非常常见的 Bug。队伍分别调速度环和位置环，单独看都没问题，合起来机器人就振荡。修复几乎总是：加快内环，或者放慢外环，直到恢复 5 倍分离。

\textbf{数值例子。}假设速度环闭环带宽 300~Hz，你把位置环带宽设到 200~Hz。比值只有 $300/200 = 1.5\times$。在位置环的穿越频率（200~Hz）处，速度环闭环传递函数贡献的相位滞后为：
\[
\angle T_{\text{vel}}(j \cdot 2\pi \cdot 200) \approx -\arctan\!\left(\frac{200}{300}\right) = -33.7°
\]

多出了 34\textdegree{} 的意外相位滞后。如果位置控制器设计时留了 45\textdegree{} 的相位裕度，现在只剩 $45 - 34 = 11$\textdegree{}——离失稳危险得很。位置环剧烈超调，可能振荡。如果把位置环带宽降到 60~Hz（5 倍分离），速度环只贡献 $-\arctan(60/300) = -11.3$\textdegree{} 的相位滞后。相位裕度从 45\textdegree{} 降到一个舒适的 34\textdegree{}，系统表现良好。

\subsubsection{调参顺序：从里往外}

带宽分离原则规定了严格的调参流程：

\textbf{第一步：}关闭所有外环，直接发电流指令。调电流 PI 控制器，直到电流响应又快又稳（超调 $<5\%$，调节时间 $<1$~ms）。这一步通常由电机驱动厂家做好了，所以你往往可以跳过。

\textbf{第二步：}闭合电流环。现在发速度指令。电流环是你的"快执行器"——不用管它了。调速度 PID 直到速度响应满意（超调 $<10\%$，调节时间 $<20$~ms）。使用微分作用于测量值（Derivative-on-Measurement）来避免设定值突变引起的尖峰。

\textbf{第三步：}闭合速度环。现在发位置指令。速度环是你的"中速执行器"。调位置 PD 或 PID 以获得期望的位置响应。

关键原则：\textit{闭合外环之后，绝不回头重调内环}，除非出了根本性的问题（比如换了电机导致绕组电阻变了）。如果你发现自己在调速度环来修位置环的振荡，真正的问题几乎一定是带宽分离不足。

比赛赶工时常见的错误：位置环振荡了，有人赶紧把速度环的 $K_p$ 调小来"压一压"。这让速度环变得迟钝，外面的位置环得不到足够的带宽，反而振得更厉害，于是又有人减位置 $K_p$，结果整个系统又慢又软。正确做法是检查带宽比，而不是无脑降增益。

要验证串级是否健康，同时记录每一层的设定值和测量值。画速度设定值 vs 实际速度：实际值应该以很小的延迟和超调跟踪设定值。如果速度响应慢或者震荡，先修好它\textit{再}碰位置环。只有速度环看起来利索了，才进入位置环调参。

%-----------------------------------------------------------------------------
\subsection{任务空间 vs 关节空间控制}
%-----------------------------------------------------------------------------

你的四轮麦克纳姆底盘需要同时控制前进、横移和旋转三个自由度，但你有四个电机。指令该发给谁？每个轮子独立控制，还是统一协调？这就引出了两种控制哲学。

\textbf{关节空间控制}（Joint-level Control）给每个电机配独立的控制器。电机 1 有一个 PID 跟踪它的速度设定值，电机 2 有自己的 PID，电机 3 也有自己的。它们互不通信。这很简单、好调试，在电机之间天然解耦时效果很好——轮式机器人每个轮子有独立速度设定值时就是这种情况。

\textbf{任务空间控制}（Task-level Control）在\textit{任务空间}（也叫笛卡尔空间或操作空间）里工作。控制器在对任务有意义的空间中计算误差和指令——比如全向底盘的 $(v_x, v_y, \omega)$——然后用逆运动学或雅可比矩阵映射到各关节指令。

什么时候用哪种？差速底盘用关节空间控制就行了：左轮速度和右轮速度到 $(v, \omega)$ 的映射是平凡的。但对于四个麦克纳姆轮的全向底盘，任务空间控制几乎是必须的。

\textbf{例子：全向底盘逆运动学。}假设你要底盘在体坐标系下以 $(v_x, v_y, \omega)$ 运动。四个麦轮按标准 X 型布局，每个轮子到中心的距离为 $l$（沿对角线量），各轮速度为：

\begin{equation}
\begin{bmatrix} \omega_1 \\ \omega_2 \\ \omega_3 \\ \omega_4 \end{bmatrix}
= \frac{1}{r}
\begin{bmatrix}
1 & -1 & -(l_x + l_y) \\
1 &  1 &  (l_x + l_y) \\
1 &  1 & -(l_x + l_y) \\
1 & -1 &  (l_x + l_y)
\end{bmatrix}
\begin{bmatrix} v_x \\ v_y \\ \omega \end{bmatrix}
\end{equation}

其中 $r$ 是轮半径，$l_x$ 是前后半长（前轮到中心），$l_y$ 是左右半宽（侧面到中心）。这个矩阵就是\textit{逆运动学}——把任务空间指令映射到关节空间。每个轮子接着用自己的速度 PID 跟踪计算出的 $\omega_i$。轨迹跟踪控制器完全在任务空间操作，对单个轮子浑然不觉。

任务空间控制的优势是：轨迹规划器用对任务有意义的坐标思考——"以 2~m/s 前进同时以 1~rad/s 旋转"。逆运动学搞定四个轮子如何协调配合这些麻烦细节。

注意上面的逆运动学矩阵\textit{不是方阵}——它把 3 个任务空间自由度映射到 4 个轮速。系统具有运动学冗余（执行器比自由度多）。这种冗余其实是好事：底盘可以实现任何 $(v_x, v_y, \omega)$ 组合，制造公差导致的轮速微小差异也能被优雅地吸收。对于执行器数等于自由度数的系统（比如三轮全向机器人），矩阵是方阵且双向可逆——你还可以做\textit{正运动学}，从测量轮速估计底盘速度，这对里程计很有用。

对于机械臂和多自由度机构——有俯仰和方位的炮塔，或者取样机械臂——任务空间方法就变得必不可少了。操作员或自主规划器想的是"把末端执行器移到 $(x, y, z)$"，逆运动学把这解算成关节角度。在这类系统里如果独立控制每个关节，末端执行器的轨迹会弯弯曲曲、不可预测。

%-----------------------------------------------------------------------------
\subsection{上层管理模式}
%-----------------------------------------------------------------------------

\subsubsection{自主运行的状态机}

控制回路计算力矩和速度。但\textit{谁来决定机器人什么时候该动起来？}这是上层管理器（Supervisor）的活儿——一个高于控制的逻辑层，管理机器人的运行模式。

最简单也最稳健的管理器是\textbf{有限状态机}（Finite State Machine, FSM）。一台典型的比赛机器人有这些状态：

\begin{center}
\texttt{IDLE} $\to$ \texttt{INITIALIZE} $\to$ \texttt{RUN} $\to$ \{\texttt{FAULT}, \texttt{E\_STOP}\}
\end{center}

在 \texttt{IDLE} 状态，所有执行器禁用，等待启动信号（如裁判系统指令）。在 \texttt{INITIALIZE} 状态，机器人执行自检：IMU 标定、编码器归零、CAN 总线健康检查。只有所有检查通过后，才跳转到 \texttt{RUN}，控制回路开始工作。如果检测到任何严重故障（电机过温、通信丢失、低电量），FSM 跳转到 \texttt{FAULT}，禁用执行器并记录错误。急停信号从任何状态直接跳转到 \texttt{E\_STOP}。

为什么这很重要？没有状态机的话，比赛现场常见的灾难是这样的：机器人上电，PID 回路立刻开始跑，IMU 还没标定完，"倾角"读出来是垃圾，电机直接满功率输出。有了状态机，机器人在初始化完成前\textit{进不了} \texttt{RUN}。

一个最小化的 C 实现：

\begin{verbatim}
    typedef enum { IDLE, INIT, RUN, FAULT, ESTOP } State;
    State state = IDLE;

    void supervisor_update(void) {
        switch (state) {
            case IDLE:  if (start_signal()) state = INIT;  break;
            case INIT:  if (all_checks_ok()) state = RUN;  break;
            case RUN:   if (fault_detected()) state = FAULT; break;
            case FAULT: motors_disable(); break;
            case ESTOP: motors_disable(); break;
        }
        if (estop_pressed()) state = ESTOP;
    }
\end{verbatim}

十行代码。让你的机器人在上电瞬间不至于自毁。

\subsubsection{模式切换与无扰切换}

比赛机器人经常需要在运行中切换控制模式。云台可能从手动模式（操作员摇杆输入）切换到自动模式（视觉跟踪）。底盘可能从遥控切换到自主导航。我们来想想切换瞬间会发生什么：旧控制器正在输出一个力矩，新控制器被激活时，它的积分器是空的，所以它的输出从零开始。结果？控制信号突然跳一下，云台甩一下。

修复方案叫\textbf{无扰切换}（Bumpless Transfer）。思路很简单：切换到新控制器时，初始化它的内部状态，使其第一拍输出恰好等于旧控制器的最后输出。对 PID 控制器来说，就是设置积分器：

\begin{equation}
I_{\text{new}} = u_{\text{current}} - K_{p,\text{new}} \cdot e_{\text{current}}
\end{equation}

其中 $u_{\text{current}}$ 是切换瞬间的控制输出，$e_{\text{current}}$ 是新控制器看到的当前误差。这保证了切换后第一拍 $u_{\text{new}} = K_{p,\text{new}} e + I_{\text{new}} = u_{\text{current}}$。新控制器随后平滑接管。

如果新控制器还有微分项，你还应该把"上一次测量值"变量初始化为当前测量值，避免微分产生尖峰。

在比赛中，最常见的模式切换是操作员控制与自动瞄准之间的切换。手动控制时，云台跟踪摇杆输入。当视觉系统检测到目标、操作员按下"自瞄"按钮时，控制器切换到视觉引导的位置控制。没有无扰切换的话，云台会甩到视觉控制器积分器初始值对应的位置（通常是零）。有了无扰切换，过渡是无缝的——视觉控制器从操作员控制器停下来的地方精确接手，平滑地转向目标。

%-----------------------------------------------------------------------------
\subsection{时序架构}
%-----------------------------------------------------------------------------

前面讲的回路堆叠理论，如果代码不能按时执行，全都白搭。时序是撑起整个架构的骨架。

\begin{center}
\begin{tabular}{l l l l}
\hline
\textbf{回路} & \textbf{频率} & \textbf{优先级} & \textbf{定时器/触发} \\
\hline
电流 PI & 10~kHz & 最高 & ADC 转换完成中断 \\
速度 PID & 1~kHz & 高 & TIM6（1~ms 周期） \\
位置 PID & 500~Hz & 中 & 每两个速度环 tick 执行一次 \\
EKF 更新 & 1~kHz & 中 & 与速度环同步 \\
轨迹规划 & 100~Hz & 低 & 每十个速度环 tick 执行一次 \\
上层状态机 & 100~Hz & 低 & 与轨迹规划同一个 tick \\
\hline
\end{tabular}
\end{center}

\textbf{优先级调度。}当电流环跑在你的 MCU 上（简单 H 桥驱动器）时，它是最高优先级。当使用 CAN 总线电调（C610/C620）时，电流环在电调内部运行——所以在你的 MCU 上，CAN 收发和速度环成为最高优先级任务。位置环和规划器优先级最低，因为它们变化慢，偶尔有点抖动可以容忍。

\textbf{为什么时序抖动很重要。}回想 PID 微分项：$D = K_d \cdot (y_{\text{prev}} - y) / \Delta t$。如果 $\Delta t$ 因为更高优先级中断延迟了你的回路而波动，即使测量值完全平滑，微分输出也会波动。对于跑在 1~kHz 的回路，10\% 的抖动（$\pm 0.1$~ms）会给微分带来 10\% 的噪声。所以确定性的时序不是可选的——它是正确性要求。

微控制器上有两种常见方案：

\textbf{裸机（嵌套中断）：}每个回路跑在定时器中断里。电流环中断在 NVIC 中优先级最高，速度环低一级，以此类推。硬件保证高优先级中断可以抢占低优先级中断。简单、开销小，但回路多了交互复杂时会很乱。

\textbf{RTOS（FreeRTOS、RT-Thread）：}每个回路是一个任务，分配优先级。RTOS 调度器处理抢占。对于子系统多的复杂机器人更清爽。

一个最小化的 FreeRTOS 配置：

\begin{verbatim}
    void vel_loop_task(void *arg) {
        TickType_t last = xTaskGetTickCount();
        for (;;) {
            vTaskDelayUntil(&last, pdMS_TO_TICKS(1));  // 1 kHz
            float omega = read_encoder_velocity();
            float cmd = pid_update(&vel_pid, vel_setpoint, omega);
            set_current_setpoint(cmd);
        }
    }
    // In main(): xTaskCreate(vel_loop_task, "vel", 256, NULL, 3, NULL);
\end{verbatim}

\texttt{vTaskDelayUntil} 这个函数至关重要：它保证任务以固定周期运行，不管计算花了多久（只要在周期内完成）。如果用 \texttt{vTaskDelay}，每次迭代就是"计算时间 + 1~ms"而不是精确的 1~ms，几千次迭代下来相位偏移就很明显了。

最后一个实用提醒：注意 CAN 总线通信的时序。如果主控以 1~kHz 通过 CAN 向电机电调发电流指令，而 CAN 总线上还跑着 IMU 数据、编码器数据和裁判系统数据包，总线竞争会导致不可预测的延迟。给时间关键的消息（如电流指令）分配更高的 CAN 优先级（更低的 CAN ID），并考虑错开发送时机，避免多条高速率消息在同一毫秒边界上撞车。

%-----------------------------------------------------------------------------
\subsection*{}
%-----------------------------------------------------------------------------

\vspace{0.5em}
\begin{mdframed}[linewidth=1pt, innertopmargin=10pt, innerbottommargin=10pt]
\textbf{案例分析：完整的云台控制栈——从传感器到执行器}

\vspace{0.3em}
让我们追踪一个比赛云台俯仰轴的完整数据流，从传感器到电机。这是很多冠军机器人在跑的架构，端到端地理解它正是本章的目标。

\textbf{流水线：}

\begin{enumerate}[nosep]
    \item \textbf{IMU}（BMI088）以 8~kHz 采样。原始陀螺仪读数 $\omega_{\text{raw}}$ 经过一阶低通滤波器（$f_c = 200$~Hz）去除高频振动。
    \item \textbf{EKF} 以 1~kHz 运行。它融合滤波后的陀螺仪与加速度计（已校正向心加速度），输出姿态估计 $\hat{\theta}$（俯仰角）和偏置校正后的角速度 $\hat{\omega}$。
    \item \textbf{位置 PID} 以 500~Hz 运行。计算 $e_\theta = \theta_{\text{cmd}} - \hat{\theta}$，输出速度指令 $\omega_{\text{cmd}}$。通常是 PD 控制器（或 $K_i$ 很小的 PID），因为位置被控对象本身包含一个积分器（$\dot{\theta} = \omega$）。
    \item \textbf{速度 PID} 以 1~kHz 运行。计算 $e_\omega = \omega_{\text{cmd}} - \hat{\omega}$，输出电流指令 $i_{\text{cmd}}$。使用微分作用于测量值，避免位置控制器突变 $\omega_{\text{cmd}}$ 时产生的设定值尖峰。
    \item \textbf{电流 PI + FOC} 以 $\sim$10~kHz 在 C620 电调内部运行。你的 MCU 通过 CAN 发送 16 位有符号 $i_{\text{cmd}}$；电调内部处理电流环、磁场定向控制和三相换相，并通过 CAN 回传实际电流、转子角度和转速。
    \item \textbf{电机驱动}通过空间矢量调制以 20~kHz 开关频率将占空比转换为三相电压。
\end{enumerate}

\textbf{时序总结：}
\begin{center}
\begin{tabular}{l l l}
\hline
\textbf{模块} & \textbf{频率} & \textbf{带宽} \\
\hline
电流 PI & 10~kHz & $\sim$5~kHz \\
速度 PID & 1~kHz & $\sim$200--500~Hz \\
位置 PID & 500~Hz & $\sim$50--100~Hz \\
EKF & 1~kHz & （非控制器，但更新状态） \\
\hline
\end{tabular}
\end{center}

带宽分离清晰可见：电流比速度快 $10\times$--$25\times$，速度比位置快 $4\times$--$10\times$。每个外环都可以安全地把内环近似为单位增益。

\textbf{关键设计决策：}

\textit{电流环为什么用 PI 不用 PID？}电流被控对象是一阶的（$L/R$ 动态）。PI 控制器对阶跃输入零稳态误差，最多贡献 90\textdegree{} 相位滞后——够了。加 D 项只会放大电流传感器噪声（电流传感器噪声不小），没有收益。

\textit{速度环为什么用微分作用于测量值？}速度设定值 $\omega_{\text{cmd}}$ 每次位置控制器运行时（每 2~ms）都会离散跳变。对误差做微分会在每次跳变时产生尖峰。对测量值做微分就避免了这个问题。

\textit{位置环为什么最慢？}因为位置是速度的积分，位置控制器的工作是生成\textit{平滑}的速度指令。跑得比必要的快了，它就会对姿态估计中的每一点小噪声做出反应，产生抖动的速度指令。

\textbf{跳过电流环会怎样？}有的队伍直接从速度 PID 输出 PWM 占空比，绕过电流控制。这"能用"——电机会转——但你失去了\textit{力矩线性度}。没有电流控制，你的指令（占空比）和实际力矩（$\tau = K_t i$）之间的关系取决于反电动势、供电电压、温度和负载。速度 PID 必须补偿所有这些非线性。增益裕度下降，对负载变化的鲁棒性降低，电池电压每漂一点你就得重调。有了电流环，速度 PID 看到的是一个漂亮的线性执行器："我要 2~A，得到 2~A。"内环把脏活累活包了。

\textbf{实用建议：}调试串级系统时，从里往外来。先用示波器看电流传感器上的阶跃响应验证电流环。然后验证速度环。然后位置环。如果外环出了问题，问题几乎从来不在外环本身——通常是内环没调好或者带宽分离不够。

\textbf{实际系统的数据。}在一个调好的比赛云台俯仰轴上，典型参数是：位置 PD，$K_p = 20$，$K_d = 0.5$（输出单位：deg/s 每 deg 误差）；速度 PID，$K_p = 30$，$K_i = 5$，$K_d = 0$（输出单位：mA 每 deg/s 误差）；电流 PI 由 C620 电调内部处理。位置阶跃响应（10\textdegree{} 阶跃）超调 $<5\%$，调节时间 $\sim$80~ms。速度阶跃响应（100~deg/s 阶跃）超调 $<10\%$，调节时间 $\sim$8~ms。这些数据给出的带宽比大约为 $1/(0.08) : 1/(0.008) \approx 12.5 : 125 = 1 : 10$——正好在带宽分离的甜蜜区间。
\end{mdframed}

\vspace{0.5em}

本章介绍的架构模式——串级控制、带宽分离、任务空间 vs 关节空间、状态机、无扰切换、时序纪律——不是什么高深理论，而是每一台设计良好的比赛机器人的标配工具箱。掌握它们，你的项目"集成"阶段——各个单独能跑的子系统合到一台机器人上——就会顺畅而非灾难。

下一章，我们从连续域设计转入离散时间实现：怎么把这些漂亮的传递函数变成在微控制器上跑的代码，要考虑有限字长和固定采样率。
